\chapter{Zero-Day Options and Dealer-Gamma Flows}\label{ch:s2:zero-day-options-and-dealer-gamma-flows}

In 2025, S\&P 500 options that expire the day they trade made up 59\% of all SPX option volume, by Cboe's count. The dealers who take the other side
hedge within the day, and their hedging can push the index further or hold it back. Baltussen, Da, Lammers and Martens found that the return in the
last 30 minutes before the close is predicted by the return over the rest of the day, in more than 60 futures markets over 1974 to 2020. They linked it
to the gamma hedging of option market makers and leveraged funds. The trade that follows is easy to write: follow the day's move when dealers are
short gamma, fade it when they are long. The hard part is knowing which. This chapter simulates zero-day options whose customer side is drawn each day,
estimates dealer gamma from open interest the way analysts do, and measures what that estimate costs. On the synthetic market the trade earns a Sharpe
ratio of 1.68 when the dealers' side is known and 0.40 when it is guessed. After costs, the guess loses. The build is \texttt{firm.gammaflow}.

\section{Zero-day options and who trades them}

A zero-day option (Book~1, chapter~26) is an option on its expiry day. Its gamma is concentrated near the strike and grows without bound as the close
approaches. A move of a few points at 3:45 p.m.\ changes its delta far more than the same move at 10 a.m. Whoever is short these options holds a
position whose hedge changes quickly and whose size depends on where the index is.

\begin{dated}{2026-09}{Zero-day option volume}\label{dat:s2:zero-day-options-and-dealer-gamma-flows:volume}
Cboe reported that SPX zero-day options averaged a record 2.3 million contracts a day in 2025, 59\% of all SPX option volume. In August 2025 their
share reached a record 62.4\%, about 2.4 million contracts a day, with retail traders an estimated 53\% of that volume.
\end{dated}

Who holds which side matters more than the volume. If customers buy the options, dealers are short gamma: they buy the index after it rises and sell
after it falls, and their hedging adds to the move. If customers sell them, dealers are long gamma and lean against the move. Ni, Pearson, Poteshman
and White found evidence of this channel in single stocks. Option market makers' hedge rebalancing, with no information behind it, affected stock
return volatility and the probability of large price moves.

\section{Estimating dealer gamma}

Open interest by strike is public. The side of each contract, and so who is long and who is short, is not. An analyst who wants dealer gamma must
assume it.

\begin{definition}[Gamma exposure estimate]\label{def:s2:zero-day-options-and-dealer-gamma-flows:gex}
A \emph{gamma exposure estimate}\index{gamma exposure estimate} is the dealers' total gamma computed from open interest by strike and an assumption
about which side dealers hold, usually that they are long the calls customers sell and short the puts customers buy: $\sum_K (\text{call OI}_K -
\text{put OI}_K)\,\Gamma_K$ under that convention.
\end{definition}

\begin{definition}[Gamma flip level]\label{def:s2:zero-day-options-and-dealer-gamma-flows:flip}
The \emph{gamma flip level}\index{gamma flip level} is the index level at which an estimate of the dealers' total gamma changes sign: above it
dealers are estimated long gamma, below it short.
\end{definition}

\texttt{firm.gammaflow} simulates 2,000 days. The index opens at 100, and zero-day calls and puts are listed on 61 strikes every 0.1 from 97 to 103.
Call open interest is centred above the open and put open interest below. Each is scaled by a daily random amount, with more on round strikes. Each
day the customers' net position in calls and in puts is drawn as a share of open interest between $-1$ (short) and 1 (long). On average customers
are short 0.2 of the calls and long 0.3 of the puts, with a standard deviation of 0.5 from day to day, and the dealers hold the opposite. The dealers'
options gain delta as the index moves. Dealers trade a tenth of that unhedged delta every five minutes and all that remains over the last 30 minutes,
and each trade moves the next bar's price (\cref{lst:s2:zero-day-options-and-dealer-gamma-flows:days}).

The truth has dealers short gamma at the open on 56.4\% of days. The conventional estimate (dealers long every call, short every put) calls 51.5\% of
days short. It agrees with the truth's sign on 56.0\% of days, and its correlation with the true gamma is 0.23
(\cref{fig:s2:zero-day-options-and-dealer-gamma-flows:estimate}). Knowing the customers' average sides does barely better, 57.5\% and 0.24, because
what the estimate misses is the day's own side, not the average. The conventional estimate finds a flip level between 97 and 103 on 96.0\% of the
first 200 days, with a median of 100.28. The truth has one on only 36.5\% of them: on most days the dealers' true gamma has the same sign across the
whole range.

\begin{omfigure}
\begin{tikzpicture}
\begin{axis}[width=9cm,height=6.5cm,xlabel={\small true dealer gamma at the open},ylabel={\small conventional estimate},
  tick label style={font=\footnotesize},xmin=-25,xmax=32,ymin=-25,ymax=30,grid=major,grid style={gray!25},table/col sep=comma]
\addplot[only marks,mark=*,mark size=0.8pt,omThm,opacity=0.6] table[x=true,y=convention]{figdata/strategies-2/06-zero-day-options-and-dealer-gamma-flows/estimate.csv};
\addplot[black,thin,domain=-25:30] {0};
\addplot[black,thin] coordinates {(0,-25) (0,30)};
\end{axis}
\end{tikzpicture}

\omcaption{The dealers' true total gamma at the open against the conventional estimate from open interest (dealers long every call, short every put),
first 400 of the 2,000 synthetic days, in the model's units. Points in the upper-left and lower-right quadrants are days whose sign the estimate gets
wrong: 44\% of all days. Data: \texttt{s2\_gammaflow.days}.}\label{fig:s2:zero-day-options-and-dealer-gamma-flows:estimate}
\end{omfigure}

\section{Pinning and pushing}

\begin{definition}[Pinning]\label{def:s2:zero-day-options-and-dealer-gamma-flows:pinning}
\emph{Pinning}\index{pinning} is the tendency of an underlying to close near a strike with large open interest on the options' expiry day, attributed
to the hedging of dealers who are long gamma at that strike: they sell as the price rises through it and buy as it falls.
\end{definition}

Pushing is the other regime: short-gamma dealers amplify moves. The simulation shows it clearly. On days when dealers were truly short gamma at the
open, the day's log return had a standard deviation of 1.48\%. On long-gamma days it was 0.81\%, against 0.99\% for the same shocks without any
hedging feedback.

\omterm{def:s2:zero-day-options-and-dealer-gamma-flows:pinning}{Pinning} is another matter. The share of closes within 0.05 of a round strike (a multiple of 0.5) is 21.4\% on long-gamma days, 20.4\% on short-gamma
days and 21.4\% without feedback. Chance alone gives 20\%. The model damps moves when dealers are long gamma, but nothing in it pulls the price to a
particular strike: dealers here hedge a share of their unhedged delta, which leans against the move wherever the index is. A pin needs gamma
concentrated at one strike in the last minutes and hedging that tracks it closely. Before trading a pin, check that your model, or your market,
produces one.

\section{Trading the hedging flow}

The trade follows Baltussen and co-authors: at 3:30 p.m.\ go long if the index is up since the open, short if it is down, and close at 4:00 p.m.
Its return in basis points of the index, by the dealers' gamma at the open, and a rule that follows the move on short-gamma days and fades it
on long-gamma days:

\begin{center}
\small
\begin{tabular}{@{}lccccc@{}}
\toprule
2,000 days, bp a day & short gamma (t) & long gamma (t) & rule & Sharpe & after 1 bp\\
\midrule
true side & 3.12 (3.0) & $-3.40$ ($-4.1$) & 3.24 & 1.68 & 1.16\\
conventional estimate & 1.02 (1.0) & $-0.52$ ($-0.5$) & 0.78 & 0.40 & $-0.11$\\
average sides & 1.08 (1.3) & $-1.77$ ($-1.4$) & 1.28 & 0.66 & 0.14\\
\bottomrule
\end{tabular}
\end{center}

\begin{omfigure}
\begin{tikzpicture}
\begin{axis}[width=9cm,height=5.4cm,ybar,bar width=14pt,ylabel={\small last 30 minutes (bp a day)},tick label style={font=\footnotesize},
  xmin=0.5,xmax=3.5,ymin=-4,ymax=4,xtick={1,2,3},xticklabels={true side,convention,average sides},grid=major,grid style={gray!25},
  table/col sep=comma,legend style={font=\scriptsize,at={(0.5,-0.2)},anchor=north},legend columns=2]
\addplot[fill=omThm!70,draw=omThm] table[x=x,y=short]{figdata/strategies-2/06-zero-day-options-and-dealer-gamma-flows/regimes.csv};
\addplot[fill=omDef!60,draw=omDef] table[x=x,y=long]{figdata/strategies-2/06-zero-day-options-and-dealer-gamma-flows/regimes.csv};
\legend{dealers short gamma,dealers long gamma}
\end{axis}
\end{tikzpicture}

\omcaption{The intraday momentum trade (the sign of the day's move to 3:30 p.m.\ times the last 30 minutes' return) on the 2,000 synthetic days,
split by the dealers' gamma at the open: the true side, the conventional estimate and an estimate with the customers' average sides. Data:
\texttt{s2\_gammaflow.regimes}.}\label{fig:s2:zero-day-options-and-dealer-gamma-flows:regimes}
\end{omfigure}

With the true side the regime split is sharp (\cref{fig:s2:zero-day-options-and-dealer-gamma-flows:regimes}): momentum on short-gamma days,
reversal on long-gamma days, and a rule that uses both earns 3.24 basis points a day. The conventional estimate keeps a quarter of it, 0.78 basis
points a day, a Sharpe ratio of 0.40 before costs and $-0.11$ after one basis point a round trip. Without any regime, following the day's move earns
0.27 basis points a day, a Sharpe ratio of 0.14. The mechanism is real in the simulation. Most of its value lies in information the public data do not
contain.

\section{Strategy files}

\begin{strategyfile}{Long gamma below the flip level}\label{strat:s2:zero-day-options-and-dealer-gamma-flows:flip}
\sfield{Who pays you, and why} Nobody pays; the position is protection bought where dealer hedging is estimated to amplify moves.
\sfield{Instruments and venues} Short-dated index options or futures options.
\sfield{Signal} The index below an estimated flip level.
\sfield{Sizing and execution} Small; the estimate's sign is wrong on many days.
\sfield{Costs} Option spreads; time decay when nothing happens.
\sfield{How it dies} A flip level that is an artefact of the side convention: in the chapter's simulation the convention finds one on 96\% of days and
the truth on 36.5\%.
\sfield{Horizon, capacity, infrastructure} Intraday to days; open interest by strike.
\sfield{Backtest honestly} Use open interest as it was known before the session.
\sfield{Sources} No performance figure verified.
\end{strategyfile}

\begin{strategyfile}{Pin trade near large strikes}\label{strat:s2:zero-day-options-and-dealer-gamma-flows:pin}
\sfield{Who pays you, and why} Buyers of expiring options near a strike where dealers are long gamma, if the index is held there.
\sfield{Instruments and venues} Short straddles or butterflies at the strike with the largest open interest, late on expiry day.
\sfield{Signal} Large open interest at a nearby strike and dealers estimated long gamma there.
\sfield{Sizing and execution} Small; the loss if the index breaks away is fast and large.
\sfield{Costs} Spreads on expiring options.
\sfield{How it dies} No pin: in the chapter's model closes land near round strikes as often as chance says.
\sfield{Horizon, capacity, infrastructure} The last hour of expiry day.
\sfield{Backtest honestly} Compare with the share of closes near a strike by chance.
\sfield{Sources} No performance figure verified.
\end{strategyfile}

\begin{strategyfile}{Intraday dealer-hedging momentum}\label{strat:s2:zero-day-options-and-dealer-gamma-flows:momentum}
\sfield{Who pays you, and why} Dealers and leveraged funds who must trade with the day's move near the close to rebalance their hedges.
\sfield{Instruments and venues} Index futures in the last 30 minutes.
\sfield{Signal} The sign of the day's return to 3:30 p.m., conditioned on estimated dealer gamma.
\sfield{Sizing and execution} Enter at 3:30 p.m., exit at the close.
\sfield{Costs} Two futures trades a day.
\sfield{How it dies} A wrong estimate of the dealers' side; crowding into the same close.
\sfield{Horizon, capacity, infrastructure} Thirty minutes; intraday data.
\sfield{Backtest honestly} Costs on every day; the regime estimate as it could have been made that morning.
\sfield{Sources} Baltussen, Da, Lammers and Martens (2021); this chapter: 1.68 with the true side, 0.40 with the convention, $-0.11$ after costs.
\end{strategyfile}

\section{Tutorial: who is short gamma}\label{tut:s2:zero-day-options-and-dealer-gamma-flows:tut}

\begin{tutorial}
\textbf{Goal.} Simulate zero-day options with random customer sides, estimate dealer gamma from open interest, and trade the hedging flow by
regime. \textbf{End state:} the table and the two figures.
\begin{tutsteps}
\item \textbf{The days}.
\omcode{code/firm/gammaflow/firm_gammaflow.py}{87}{118}{Open interest, customer sides, dealer hedging with lags, and the estimate under the convention.}\label{lst:s2:zero-day-options-and-dealer-gamma-flows:days}
\item \textbf{The trade by regime}.
\omcode{code/strategies-2/06-zero-day-options-and-dealer-gamma-flows/python/s2_gammaflow.py}{49}{71}{The last-30-minute trade, split by true and estimated regime.}\label{lst:s2:zero-day-options-and-dealer-gamma-flows:trade}
\item \textbf{Run} \texttt{estimates()}, \texttt{regimes()}, \texttt{day\_vol()}, \texttt{pinning()}, \texttt{flips()} and
\texttt{fig\_gammaflow.py}.
\end{tutsteps}
\textbf{What to change next.} Let dealers hedge each option's delta at every bar and see whether \omterm{def:s2:zero-day-options-and-dealer-gamma-flows:pinning}{pinning} appears; give the analyst a noisy
signal of the day's customer side and measure the trade as the signal improves; add leveraged funds' close rebalancing.
\end{tutorial}

\section{Build: gamma flow}\label{bld:s2:zero-day-options-and-dealer-gamma-flows:gammaflow}

\begin{build}
\bfield{Purpose} Zero-day open interest by strike, customer sides, dealer gamma and its conventional estimate, flip levels, and an intraday
simulator in which dealer hedging moves prices.
\bfield{Interface} \texttt{FlowConfig(\dots)}, \texttt{strikes(cfg)}, \texttt{open\_interest(cfg, centre)}, \texttt{gamma(S, K, tau, vol)},
\texttt{dealer\_gamma(S, tau, pos\_c, pos\_p, cfg)}, \texttt{flip\_level(pos\_c, pos\_p, cfg, tau)}, \texttt{simulate\_days(cfg)}.
\bfield{Rules} Dealers hold the opposite of customers; hedge trades move the next bar; the estimate sees open interest only.
\bfield{Acceptance tests} \texttt{code/firm/gammaflow/tests/}: gamma peaks at the strike and grows near expiry; no impact gives the free paths; a flip
level between long calls and short puts.
\bfield{Stretch} Per-option hedging; leveraged-fund flows; a noisy side signal.
\end{build}

\begin{omsources}
\item Cboe, trading volume reports and Cboe Insights on SPX zero-day options, 2025.
\item G.~Baltussen, Z.~Da, S.~Lammers and M.~Martens, ``Hedging demand and market intraday momentum'', \emph{Journal of Financial Economics}
142(1), 2021.
\item S.~X.~Ni, N.~D.~Pearson, A.~M.~Poteshman and J.~White, ``Does option trading have a pervasive impact on underlying stock prices?'',
\emph{Review of Financial Studies} 34(4), 2021.
\end{omsources}

\section{Exercises}

\begin{exercise}[$\star$]\label{exo:s2:zero-day-options-and-dealer-gamma-flows:1}
Dealers are short gamma of 5 index units per index point. The index rises 2 points. What do they trade to stay hedged?
\end{exercise}

\begin{exercise}[$\star$]\label{exo:s2:zero-day-options-and-dealer-gamma-flows:2}
Why does a zero-day option's gamma near the strike grow as the close approaches?
\end{exercise}

\begin{exercise}[$\star$]\label{exo:s2:zero-day-options-and-dealer-gamma-flows:3}
Chance alone puts 20\% of closes within 0.05 of a multiple of 0.5. Why?
\end{exercise}

\begin{exercise}[$\star\star$]\label{exo:s2:zero-day-options-and-dealer-gamma-flows:4}
Why does the conventional estimate find a flip level on almost every day?
\end{exercise}

\begin{exercise}[$\star\star$]\label{exo:s2:zero-day-options-and-dealer-gamma-flows:5}
Knowing the customers' average sides improves the estimate only a little. Why?
\end{exercise}

\begin{exercise}[$\star\star$]\label{exo:s2:zero-day-options-and-dealer-gamma-flows:6}
Why does the day's volatility differ between short-gamma and long-gamma days, and what does Ni and co-authors' result say about it?
\end{exercise}

\begin{exercise}[$\star\star\star$]\label{exo:s2:zero-day-options-and-dealer-gamma-flows:7}
\emph{Coding.} Set \texttt{FlowConfig(impact=0.0)} and rerun the trade. What should it earn, and why?
\end{exercise}

\begin{exercise}[$\star\star\star$]\label{exo:s2:zero-day-options-and-dealer-gamma-flows:8}
\emph{Find the flaw.} ``Our gamma estimate says dealers are short today, so the last 30 minutes will follow the morning's move.''
\end{exercise}

\section{Problem: Who Is Short Gamma}

\begin{problem}[{Weekend problem --- trading dealers' hedges}]\label{pb:s2:zero-day-options-and-dealer-gamma-flows:1}
The chapter's simulated days, Cboe's figures and the public research.

\medskip\noindent\textbf{Part I --- The options.}
\begin{enumerate}
\item What is a zero-day option, and how does its gamma behave through the day?
\item Give Cboe's volume figures for 2025.
\item How do short-gamma and long-gamma dealers hedge?
\item What did Ni, Pearson, Poteshman and White find?
\end{enumerate}

\medskip\noindent\textbf{Part II --- The estimate.}
\begin{enumerate}[resume]
\item Define a \omterm{def:s2:zero-day-options-and-dealer-gamma-flows:gex}{gamma exposure estimate} and the flip level.
\item What convention do analysts use, and what does it assume?
\item How often does the estimate get the sign right, and how does it correlate with the truth?
\item How often does each find a flip level?
\end{enumerate}

\medskip\noindent\textbf{Part III --- The market.}
\begin{enumerate}[resume]
\item How do dealers hedge in the model, and how does their hedging move prices?
\item Give the day's volatility in each regime.
\item Define \omterm{def:s2:zero-day-options-and-dealer-gamma-flows:pinning}{pinning}. Does the model pin?
\item What would a model need to produce a pin?
\end{enumerate}

\medskip\noindent\textbf{Part IV --- The verdict.}
\begin{enumerate}[resume]
\item State the \emph{named result}: the error of the gamma estimate when the dealers' side is guessed, and the hedging-flow trade's return by
regime.
\item What did Baltussen and co-authors find?
\item What does the trade earn without a regime?
\item What survives costs?
\item Where would better information about sides come from?
\item How would you backtest the trade honestly?
\item Which strategy file depends most on the flip level?
\item In one sentence: what is the trade really betting on?
\end{enumerate}
\end{problem}

\section{Interview questions}

\begin{interviewq}[$\star$ \iqroles{trader}]\label{iq:s2:zero-day-options-and-dealer-gamma-flows:1}
What does it mean for dealers to be short gamma, and why would it matter for the index?
\end{interviewq}

\begin{interviewq}[$\star\star$ \iqroles{researcher}]\label{iq:s2:zero-day-options-and-dealer-gamma-flows:2}
How would you estimate dealer gamma, and how would you know whether the estimate is any good?
\end{interviewq}

\begin{interviewq}[$\star\star$ \iqroles{trader}]\label{iq:s2:zero-day-options-and-dealer-gamma-flows:3}
The index is near a strike with huge open interest an hour before expiry. What might happen, and how would you trade it?
\end{interviewq}

\begin{interviewq}[$\star\star$ \iqroles{risk}]\label{iq:s2:zero-day-options-and-dealer-gamma-flows:4}
A desk sells zero-day options every day. What is its risk, and how do you limit it?
\end{interviewq}

\begin{interviewq}[$\star\star$ \iqroles{developer}]\label{iq:s2:zero-day-options-and-dealer-gamma-flows:5}
Design a system that computes a \omterm{def:s2:zero-day-options-and-dealer-gamma-flows:gex}{gamma exposure estimate} every minute from open interest and quotes.
\end{interviewq}

\begin{interviewq}[$\star\star\star$ \iqroles{researcher}]\label{iq:s2:zero-day-options-and-dealer-gamma-flows:6}
Show that the gamma of an at-the-money option is $\varphi(\tfrac12\sigma\sqrt\tau)/(S\sigma\sqrt\tau)$, and find how it scales as $\tau \to 0$.
\end{interviewq}
